\cleardoublepage
\thispagestyle{empty}

\begin{center}
  {\bf \Huge Ringraziamenti}
\end{center}
\vspace{2cm}
Ora che sono arrivato alla fine di questo percorso, sento il bisogno di dedicare qualche riga a tutte quelle persone che sono state necessarie affinché la stesura di questa tesi fosse possibile e, più in generale, affinché io potessi raggiungere questo risultato.
\newline
\newline
I primi ringraziamenti vanno al mio relatore, Alberto Montresor, per il prezioso lavoro di revisione di questa tesi, e al mio co-relatore, Quintino Francesco Lotito, che mi ha seguito instancabilmente durante tutto il mio percorso di tirocinio interno.
Desidero ringraziarli non solo per il supporto tecnico e professionale che mi hanno offerto, ma soprattutto perché entrambi rappresentano per me un modello da seguire, sia dal punto di vista accademico e professionale, sia da quello personale.
Poter essere stato guidato da loro è stato dunque per me un vero e proprio privilegio ed onore.
\newline
\newline
Un ringraziamento speciale va ai miei genitori, rispettivamente co-co-relatore e co-co-relatrice per il supporto fornito durante lo sviluppo di questo lavoro.
Grazie al loro sostegno sono riuscito a finalizzare questo percorso sentendomi sempre supportato e sostenuto anche nei momenti più incerti e difficili di questo percorso.
\newline
\newline
Infine, il ringraziamento più grande va a tutte le persone che ho incontrato e che mi hanno cambiato durante questo percorso a Trento. Citarle una per una risulterebbe difficile e al contempo riduttivo; sento però il bisogno di diversificare un minimo i ringraziamenti, poiché ogni gruppo che ho frequentato e che mi ha sostenuto ha avuto un valore diverso all'interno della mia esperienza universitaria.
Vorrei prima di tutto ringraziare i gruppi che si sono formati attorno a me all'interno dello studentato Nest durante la mia permanenza a Trento.
Il gruppo del primo anno e quello del terzo anno sono stati entrambi fondamentali nello sviluppo del mio percorso accademico e personale: anche se hanno contribuito in modo diverso, sono stati e continuano a essere il motivo per cui ho deciso di proseguire il mio percorso accademico.
Vorrei inoltre ringraziare i miei compagni di corso: attraverso il loro supporto sono riuscito a instaurare un sano rapporto di competitività, che mi ha portato a un miglioramento personale, e al tempo stesso al ritrovamento di un gruppo di persone che condivide i miei stessi interessi e le mie stesse passioni.
Infine ritengo necessario un ringraziamento personale a quelle poche persone che mi hanno accompagnato durante tutti questi tre lunghissimi anni, a voi va il mio grazie più sincero dato il supporto che mi avete dimostrato durante l'interezza di questi tre anni, grazie per essere stati la mia famiglia a 400Km da casa.